\documentclass[a4paper]{article}
\usepackage{amsmath}
\usepackage{amsfonts}

\begin{document}

\noindent \large Auteur: Dina Haajou \\
\noindent \large Studierichting: Informatica \\
\noindent \large Oefeningenreeks 6A nr. 4.1 \\

\medskip

\normalsize

$x_9 = 3$ en $x_3 = 27$ \\

We gebruiken de formule voor een rekenkundige rij: $x_n = x_1 + (n - 1)v$ \\

Toegepast op $x_9$: \\

$3 = x_1 + 8v$ \\

En op $x_3$: \\

$27 = x_1 + 2v$ \\

We trekken de tweede vergelijking af van de eerste: \\

$3 - 27 = (x_1 + 8v) - (x_1 + 2v)$ \\

$-24 = 6v$ \\

$v = -4$ \\

We vullen $v$ terug in in $x_3 = x_1 + 2v$: \\

$27 = x_1 + 2(-4) = x_1 - 8$ \\

$x_1 = 35$ \\

Antwoord: $x_1 = 35$

\begin{thebibliography}{999}
% \bibitem{}
\end{thebibliography}

\end{document}
